\section{Etapa 6 --- Auto-calibraci\'on, DMA y \texttt{superMaquina}}
\label{sec:et6}
\periodo{14 de junio --- 2 de julio de 2026 \quad$\cdot$\quad 11 d\'ias, 17 \emph{commits}}

La etapa que produjo el resultado publicable del proyecto. Tiene tres l\'ineas
que avanzan entrelazadas: cerrar el lazo de calibraci\'on del \emph{front-end},
sacar la adquisici\'on de la CPU, y finalmente sacar la \emph{decisi\'on} de
adquirir tambi\'en de la CPU.

\subsection{Por qu\'e hace falta calibrar}

La cadena anal\'ogica tiene cuatro etapas en cascada ---PGA, pasa-banda (BP),
sumador (SUM) y pasa-bajos (LP)--- y cada una aporta su propio \emph{offset} de
continua: tensiones de \emph{offset} de los operacionales, corrientes de
polarizaci\'on, desapareamiento de resistencias. El problema es que
\emph{el \emph{offset} se amplifica junto con la se\~nal}, y la compensaci\'on
del ge\'ofono exige ganancia alta ---proporcional a $(\zeta_1-\zeta_0)$, con
$\zeta_0 = 0{,}25$ y $\zeta_1$ de varios cientos: tres \'ordenes de magnitud.

\dato{Dimensi\'on del problema, medida: el \emph{offset} con referencia fija a
la salida de la etapa sumadora lleg\'o a \textbf{\SI{21.8}{\percent}} del fondo
de escala, y a la salida del LP ---la que alimenta al ADC--- hasta
\textbf{\SI{45.8}{\percent}}. Es decir, casi la mitad del rango de entrada
consumido \emph{antes de aplicar se\~nal s\'ismica alguna}.}

\subsection{Los tres algoritmos que se probaron}

\begin{center}\small
\begin{tabularx}{\textwidth}{@{}l X l@{}}
\toprule
\textbf{Enfoque} & \textbf{Qu\'e pas\'o} & \textbf{Destino} \\
\midrule
Bisecci\'on binaria sobre c\'odigos VDAC
& Primer intento (14 jun), solo para GEO. Funciona pero es ciego: no aprovecha que la planta es conocida, y cada etapa necesitaba su propia versi\'on
& \textbf{Descartado} \\
\addlinespace[2pt]
Servo lento
& Convergencia por integraci\'on pura
& \textbf{Descartado} \\
\addlinespace[2pt]
PID
& Se prob\'o y se abandon\'o. El t\'ermino derivativo amplifica el ruido de medici\'on del ADC sin aportar nada en una planta cuasi-est\'atica ---en la bit\'acora qued\'o descrito como \emph{``overfitting al ruido''}
& \textbf{Descartado el D} \\
\addlinespace[2pt]
\textbf{PI con \emph{feedforward} f\'isico}
& Un \'unico controlador para las cuatro etapas de GEO y las dos de HAMMER
& \textbf{Se qued\'o} \\
\bottomrule
\end{tabularx}
\end{center}

\seleccionado{\textbf{El 23 de junio los dos algoritmos separados (GEO y
HAMMER) convergen en uno solo: PI puro con \emph{feedforward}}, detector de
\emph{lock} portado de una simulaci\'on de Simulink de referencia. Las
implementaciones de bisecci\'on y servo no se borraron de inmediato: quedaron
compiladas fuera con \texttt{\#if 0} como referencia, y se eliminaron del
archivo activo reci\'en el 26 de junio, una vez validado el reemplazo. Ese
orden ---validar antes de borrar--- vale la pena registrarlo.}

\begin{figure}[!ht]
\centering
\includegraphics[width=0.94\textwidth]{psoc_foreground_calibration_flow.drawio_horizontal.png}
\caption{Flujo de la calibraci\'on \emph{foreground} en el PSoC: las etapas se
calibran en orden, cada una hasta entrar en su banda muerta, antes de pasar a
la siguiente. Al terminar, el resultado se persiste en EEPROM indexado por
c\'odigo de ganancia.}
\label{fig:calib-flow}
\end{figure}

\begin{figure}[!ht]
\centering
\includegraphics[width=0.80\textwidth]{calibration_fsm_simplified.png}
\caption{M\'aquina de estados de la calibraci\'on, versi\'on simplificada. La
versi\'on detallada est\'a en \texttt{docs/calibration\_fsm/}.}
\label{fig:calib-fsm}
\end{figure}

\subsection{Los detalles finos del controlador}

Tres afinaciones sucesivas, cada una corrigiendo un modo de falla real:

\begin{enumerate}[itemsep=0.3em]
\item \textbf{Escala en c\'odigos de DAC, no en volts} (24 jun). El actuador es
un DAC; expresar el lazo en su unidad natural evita un factor de conversi\'on
que dependía de la ganancia de etapa.

\item \textbf{Banda muerta proporcional a la ganancia de etapa} (24 jun). No
tiene sentido exigir la misma precisi\'on absoluta a una etapa de ganancia
unitaria que a una de ganancia 32: en la primera, insistir solo agrega ciclos
de conmutaci\'on del DAC sin mejorar nada aguas abajo.

\item \textbf{Anti-\emph{windup} direccional} (25 jun). Este es el m\'as sutil.
El integrador no se congela ante \emph{cualquier} saturaci\'on contra el riel
del DAC, sino \'unicamente cuando seguir integrando \emph{en ese sentido}
empeorar\'ia la saturaci\'on. Congelarlo en ambos sentidos deja al lazo incapaz
de salir del riel. Adem\'as hubo que manejar correctamente el \textbf{signo de
la ganancia} por etapa: \texttt{HAMMER\_PGA} es inversora y din\'amica,
\texttt{HAMMER\_LP} es fija y no inversora.
\end{enumerate}

\dato{Primeras convergencias documentadas en hardware real:
\texttt{HAMMER\_PGA ok=1 dac=65}, \texttt{HAMMER\_LP ok=1 dac=57} (23 jun); al
d\'ia siguiente \texttt{HAMMER\_LP} converge con \texttt{err\_mV=0} exacto
contra un objetivo de \SI{1024}{\milli\volt}. El 26 de junio el PI corre por
primera vez sobre las cuatro etapas de GEO.}

\subsection{Dos bugs que valen su propia entrada}

\dato{\textbf{El filtro cargaba los coeficientes equivocados} (24 jun). Se
descubri\'o que el filtro de adquisici\'on llevaba \emph{meses} cargando
coeficientes de pasa-banda donde deb\'ia haber un pasa-bajos. No se manifestaba
como falla evidente: el sistema funcionaba, simplemente filtraba otra cosa.}

\dato{\textbf{Un pin sin soldar} (26 jun). Un problema de calibraci\'on que se
persigui\'o como bug de firmware result\'o ser hardware. Registrado porque es la
clase de error que consume horas de depuraci\'on en el dominio equivocado.}

\subsection{Persistencia: la calibraci\'on sobrevive al reset}

El 17 de junio la calibraci\'on empez\'o a guardarse en EEPROM con CRC-16. El
1\textsuperscript{o} de julio se rehizo el esquema para guardar
\textbf{un \emph{slot} por c\'odigo de ganancia del PGA}, porque el punto de
operaci\'on \'optimo depende de la ganancia seleccionada.

\seleccionado{Consecuencia operativa: con calibraci\'on almacenada v\'alida, el
arranque solo ejecuta los intervalos de establecimiento por etapa
(2560 muestras, \textasciitilde{}\SI{1.0}{\second} a \SI{2604}{\hertz}) en vez
de la calibraci\'on completa. La diferencia entre esperar un segundo y esperar
un minuto es la diferencia entre un instrumento usable en campo y uno que no.}

\subsection{La adquisici\'on sale de la CPU}

Tres pasos en tres d\'ias:

\begin{center}\small
\begin{tabularx}{\textwidth}{@{}l l X@{}}
\toprule
\textbf{Fecha} & \textbf{Cambio} & \textbf{Consecuencia} \\
\midrule
17 jun & ISR cl\'asica $\to$ cadena DMA con \texttt{Filter} FIR de hardware &
El ARM deja de intervenir en el camino de la muestra. El FIR corre en el bloque dedicado, no en software \\
\addlinespace[2pt]
18 jun & Multiplexor DMA de cuatro canales (crudo/filtrado $\times$ simple/lote) gobernado por \texttt{Reg\_Select} &
Se cablea el promediado por hardware pero se lo deja inerte: la ruta existe antes de usarse \\
\addlinespace[2pt]
19 jun & Se activa el canal ``Lote'' &
La calibraci\'on deja de promediar muestra a muestra en el ISR y consume ventanas ya sumadas por hardware, con verificaci\'on \emph{en tiempo de compilaci\'on} de que el tama\~no del lote DMA coincide con la ventana de calibraci\'on \\
\bottomrule
\end{tabularx}
\end{center}

\subsection{\texttt{superMaquina}: la captura se vuelve hardware (2 de julio)}

El paso final. La l\'ogica de captura ---que hasta ese momento viv\'ia en C,
como ISRs sobre \texttt{Reg\_Select}--- se reescribi\'o como una \textbf{m\'aquina
de estados en Verilog dentro de un bloque UDB} del PSoC.

\texttt{superMaquina} decide, en hardware:

\begin{itemize}[itemsep=0.2em]
\item qu\'e DRQ de DMA se deja pasar;
\item cu\'ando arranca y cu\'ando termina una captura;
\item qu\'e interrupciones llegan efectivamente al ARM.
\end{itemize}

\seleccionado{\textbf{Por qu\'e importa tanto para MASW.} Mientras la captura la
gobierna software, el instante de inicio y el intervalo entre muestras heredan
el \emph{jitter} de las interrupciones ---latencia variable, prioridades,
secciones cr\'iticas. En hardware, el temporizado es determinista por
construcci\'on. Como el m\'etodo mide velocidad de fase, y la fase se mide contra
el tiempo, \emph{determinismo temporal es exactitud de medici\'on}.}

\begin{figure}[!ht]
\centering
\includegraphics[width=\textwidth]{DigitalPathFinal_trim.png}
\caption{Camino digital dentro del PSoC tras esta etapa: ADC $\Delta\Sigma$
$\to$ \texttt{Filter} FIR de hardware $\to$ cadena DMA con canales
crudo/filtrado y simple/lote, todo gobernado por \texttt{superMaquina}. El ARM
queda fuera del camino de la muestra.}
\label{fig:digital}
\end{figure}

En paralelo, tres \emph{timers} fijos del PSoC reemplazaron el \emph{tick} de
sistema de \SI{10}{\milli\second} (ping de arranque/reposo, parpadeo de LED,
\emph{watchdog} y progreso de calibraci\'on), y el ESP32 esclavo gan\'o un juego
completo de comandos por USB para operar el \emph{pipeline} paso a paso desde
la terminal, sin depender de la web: \texttt{probe}, \texttt{status},
\texttt{stream}, \texttt{debugpsoc}, \texttt{pre}, \texttt{sync},
\texttt{startnow}, \texttt{cap}, \texttt{clear}, \texttt{stop}.

\seleccionado{Ese juego de comandos result\'o ser la herramienta de diagn\'ostico
m\'as usada del proyecto: permite reproducir en aislamiento cualquier paso de
la secuencia de captura, que es exactamente lo que hizo falta para encontrar
los bugs cr\'iticos de la Etapa~8.}

\pendiente{Pol\'itica de eventos que qued\'o fijada y conviene no romper: el
\texttt{TC} se atiende por \emph{polling} de \texttt{STATUS\_TC} y \textbf{no}
por interrupci\'on; la ISR maneja \'unicamente muestras; y \texttt{CLEAR\_FLAGS}
est\'a prohibido durante \texttt{SAMPLING}.}

\subsection{Resultado medido de la auto-calibraci\'on}

\begin{figure}[!ht]
\centering
\includegraphics[width=0.92\textwidth]{autocalibration_stage_outputs_5_to_20s.png}
\caption{Trayectorias de calibraci\'on de las cuatro etapas anal\'ogicas
capturadas con osciloscopio. La correcci\'on visible abarca unos
\SI{7}{\second}. Cada etapa se calibra en orden y su salida queda dentro de la
banda muerta antes de pasar a la siguiente.}
\label{fig:calib-traj}
\end{figure}

Sobre \textbf{dos placas ensambladas independientemente} y \textbf{ocho
configuraciones placa--ganancia}, con $N=10$ reinicios apareados por
configuraci\'on:

\begin{center}\small
\begin{tabularx}{\textwidth}{@{}l X@{}}
\toprule
\emph{Offset} residual a la salida LP & de \SIrange{2.9}{45.8}{\percent} del fondo de escala a \SIrange{0.4}{3.9}{\percent} \\
Observaciones individuales calibradas & entre \SIrange{0.1}{5.2}{\percent} \\
Prueba de signos de una cola, etapa LP & rechaza la hip\'otesis nula en \textbf{las ocho} configuraciones (10 de 10 reinicios, $p\approx10^{-3}$ exacto); \textbf{sobrevive a la correcci\'on de Bonferroni} \\
Ajuste del controlador & \textbf{\'unica configuraci\'on} para todas las placas y ganancias, sin reajuste por placa ni por ganancia \\
Tiempo & correcci\'on visible \textasciitilde{}\SI{7}{\second}; completitud informal \textasciitilde{}\SI{10}{\second}; ventanas l\'imite configuradas \SI{57.6}{\second} \\
\bottomrule
\end{tabularx}
\end{center}

\dato{\textbf{El \'unico vencimiento de tiempo l\'imite observado en banco} fue
el de una placa alimentada mientras su ge\'ofono estaba en movimiento: el
transitorio de alterna viola la hip\'otesis cuasi-est\'atica que la calibraci\'on
\emph{foreground} requiere. Es una limitaci\'on entendida, no una falla.}

\subsection{Balance de la etapa}

\begin{center}\small
\begin{tabularx}{\textwidth}{@{}l X@{}}
\toprule
\textbf{Se qued\'o} & Controlador PI \'unico con \emph{feedforward}, banda muerta proporcional y anti-\emph{windup} direccional; persistencia en EEPROM con CRC-16 indexada por c\'odigo de ganancia; adquisici\'on por DMA con FIR en hardware; promediado por lote; \texttt{superMaquina} en UDB; comandos USB de diagn\'ostico \\
\addlinespace[3pt]
\textbf{Se descart\'o} & Bisecci\'on binaria; servo lento; t\'ermino derivativo del PID; \texttt{Reg\_Select} como gobierno de la captura; \emph{tick} de sistema de \SI{10}{\milli\second} \\
\addlinespace[3pt]
\textbf{Qued\'o abierto} & Deriva de largo plazo y apareamiento entre canales del arreglo completo, sin caracterizar. La evidencia obtenida concierne al \emph{punto de operaci\'on} de la cadena, no al desempe\~no s\'ismico extremo a extremo \\
\bottomrule
\end{tabularx}
\end{center}
